\documentclass[10pt,a4paper]{amsart}
\usepackage{amsmath,amssymb,amsthm,graphicx}
\usepackage{verbatim}
\newenvironment{reference}{\comment}{\endcomment}
%\newenvironment{reference}{}{}
\newenvironment{slogan}{\comment}{\endcomment}
\newenvironment{history}{\comment}{\endcomment}

% For commutative diagrams we use Xy-pic
\usepackage[all]{xy}

% We use 2cell for 2-commutative diagrams.
\xyoption{2cell}
% Theorem environments.
%
\theoremstyle{plain}
\newtheorem{theorem}[subsection]{Theorem}
\newtheorem{proposition}[subsection]{Proposition}
\newtheorem{lemma}[subsection]{Lemma}

\theoremstyle{definition}
\newtheorem{definition}[subsection]{Definition}
\newtheorem{example}[subsection]{Example}
\newtheorem{exercise}[subsection]{Exercise}
\newtheorem{situation}[subsection]{Situation}

\theoremstyle{remark}
\newtheorem{remark}[subsection]{Remark}
\newtheorem{remarks}[subsection]{Remarks}

\numberwithin{equation}{subsection}

% Macros
%
\def\lim{\mathop{\mathrm{lim}}\nolimits}
\def\colim{\mathop{\mathrm{colim}}\nolimits}
\def\Spec{\mathop{\mathrm{Spec}}}
\def\Hom{\mathop{\mathrm{Hom}}\nolimits}
\def\Ext{\mathop{\mathrm{Ext}}\nolimits}
\def\SheafHom{\mathop{\mathcal{H}\!\mathit{om}}\nolimits}
\def\SheafExt{\mathop{\mathcal{E}\!\mathit{xt}}\nolimits}
\def\Sch{\mathit{Sch}}
\def\Mor{\mathop{\mathrm{Mor}}\nolimits}
\def\Ob{\mathop{\mathrm{Ob}}\nolimits}
\def\Sh{\mathop{\mathit{Sh}}\nolimits}
\def\NL{\mathop{N\!L}\nolimits}
\def\CH{\mathop{\mathrm{CH}}\nolimits}
\def\proetale{{pro\text{-}\acute{e}tale}}
\def\etale{{\acute{e}tale}}
\def\QCoh{\mathit{QCoh}}
\def\Ker{\mathop{\mathrm{Ker}}}
\def\Im{\mathop{\mathrm{Im}}}
\def\Coker{\mathop{\mathrm{Coker}}}
\def\Coim{\mathop{\mathrm{Coim}}}

% Boxtimes
%
\DeclareMathSymbol{\boxtimes}{\mathbin}{AMSa}{"02}

%
% Macros for moduli stacks/spaces
%
\def\QCohstack{\mathcal{QC}\!\mathit{oh}}
\def\Cohstack{\mathcal{C}\!\mathit{oh}}
\def\Spacesstack{\mathcal{S}\!\mathit{paces}}
\def\Quotfunctor{\mathrm{Quot}}
\def\Hilbfunctor{\mathrm{Hilb}}
\def\Curvesstack{\mathcal{C}\!\mathit{urves}}
\def\Polarizedstack{\mathcal{P}\!\mathit{olarized}}
\def\Complexesstack{\mathcal{C}\!\mathit{omplexes}}
% \Pic is the operator that assigns to X its picard group, usage \Pic(X)
% \Picardstack_{X/B} denotes the Picard stack of X over B
% \Picardfunctor_{X/B} denotes the Picard functor of X over B
\def\Pic{\mathop{\mathrm{Pic}}\nolimits}
\def\Picardstack{\mathcal{P}\!\mathit{ic}}
\def\Picardfunctor{\mathrm{Pic}}
\def\Deformationcategory{\mathcal{D}\!\mathit{ef}}
\usepackage{cleveref}
\crefname{originallemma}{original lemma}{original lemmas}
\Crefname{originallemma}{Original lemma}{Original lemmas}
\makeatletter
\newlabel{lemma-etale-morphism-topoi}{{03LP}{source}{}{sourcepointer-1}{}}
\newlabel{lemma-etale-morphism-topoi@cref}{{[originallemma][0][]03LP}{[1][0][]source}{}{}{}}
\makeatother
\crefname{originalreference}{original reference}{original references}
\Crefname{originalreference}{Original reference}{Original references}
\makeatletter
\newlabel{etale-cohomology-lemma-describe-pullback}{{09XN}{source}{}{sourcepointer-2}{}}
\newlabel{etale-cohomology-lemma-describe-pullback@cref}{{[originalreference][0][]09XN}{[1][0][]source}{}{}{}}
\makeatother
\crefname{originalreference}{original reference}{original references}
\Crefname{originalreference}{Original reference}{Original references}
\makeatletter
\newlabel{sites-lemma-sheaf-coequalizer-representable}{{00WS}{source}{}{sourcepointer-3}{}}
\newlabel{sites-lemma-sheaf-coequalizer-representable@cref}{{[originalreference][0][]00WS}{[1][0][]source}{}{}{}}
\makeatother
\crefname{originallemma}{original lemma}{original lemmas}
\Crefname{originallemma}{Original lemma}{Original lemmas}
\makeatletter
\newlabel{lemma-support-subsheaf-final}{{04K8}{source}{}{sourcepointer-4}{}}
\newlabel{lemma-support-subsheaf-final@cref}{{[originallemma][0][]04K8}{[1][0][]source}{}{}{}}
\makeatother
\crefname{originalreference}{original reference}{original references}
\Crefname{originalreference}{Original reference}{Original references}
\makeatletter
\newlabel{spaces-theorem-presentation}{{02WW}{source}{}{sourcepointer-5}{}}
\newlabel{spaces-theorem-presentation@cref}{{[originalreference][0][]02WW}{[1][0][]source}{}{}{}}
\makeatother
\crefname{originalreference}{original reference}{original references}
\Crefname{originalreference}{Original reference}{Original references}
\makeatletter
\newlabel{sites-lemma-describe-j-shriek-representable}{{03HU}{source}{}{sourcepointer-6}{}}
\newlabel{sites-lemma-describe-j-shriek-representable@cref}{{[originalreference][0][]03HU}{[1][0][]source}{}{}{}}
\makeatother
\title{A sheaf of sets of bounded size on the small etale site is the sheaf of sections of an etale space over Y}
\author{The Stacks Project Authors}
\date{stacks-project:spaces-properties.tex:local-snapshot-20261012; GNU FDL 1.2 or later}
\begin{document}
\maketitle
\noindent\textit{Editorial note.} All delivered mathematical text is the author's own, copied verbatim from the pinned Stacks Project chapter file spaces-properties.tex; no statement or proof step was rewritten, reordered or removed.\par\medskip
\noindent\textit{Editorial note.} Span recomputation by label pairing (LANE\_BRIEF 10.2): in this chapter the author's proof environment follows the statement environment as a standalone block; the spec covers the complete environment pair (statement span, proof span) as recorded in the pool candidate.\par\medskip
\noindent\textit{Editorial note.} Occupation scan: combined claim\_key + statement\_anchor + generic-line scan over all registration files (active and withdrawn) found 13 S10-geom published cards for spaces-properties.tex (prop-occupied.json, 2026-10-14); those 13 labels are excluded from this batch by design. For the 77 free labels: 0 occupied. The registry mentions field (7557 files) and focused-seeds (645) both point to the same S10-geom cards (dual-source check), no other occupation.\par\medskip
\ifdefined\hypertarget\hypertarget{sourcepointer-1}{}\fi
\noindent\textit{Source reference.} Original lemma 03LP: sources/stacks-chow/tags.txt, source lines 4721--4721. Cross-chapter standard prerequisite cited by the author inside the counted span; it is used as a pointer and its own proof is not part of this item.\par\medskip
\ifdefined\hypertarget\hypertarget{sourcepointer-2}{}\fi
\noindent\textit{Source reference.} Original reference 09XN: sources/stacks-chow/tags.txt, source lines 12536--12536. Cross-chapter standard prerequisite cited by the author inside the counted span; it is used as a pointer and its own proof is not part of this item.\par\medskip
\ifdefined\hypertarget\hypertarget{sourcepointer-3}{}\fi
\noindent\textit{Source reference.} Original reference 00WS: sources/stacks-chow/tags.txt, source lines 1194--1194. Cross-chapter standard prerequisite cited by the author inside the counted span; it is used as a pointer and its own proof is not part of this item.\par\medskip
\ifdefined\hypertarget\hypertarget{sourcepointer-4}{}\fi
\noindent\textit{Source reference.} Original lemma 04K8: sources/stacks-chow/tags.txt, source lines 5964--5964. Cross-chapter standard prerequisite cited by the author inside the counted span; it is used as a pointer and its own proof is not part of this item.\par\medskip
\ifdefined\hypertarget\hypertarget{sourcepointer-5}{}\fi
\noindent\textit{Source reference.} Original reference 02WW: sources/stacks-chow/tags.txt, source lines 3819--3819. Cross-chapter standard prerequisite cited by the author inside the counted span; it is used as a pointer and its own proof is not part of this item.\par\medskip
\ifdefined\hypertarget\hypertarget{sourcepointer-6}{}\fi
\noindent\textit{Source reference.} Original reference 03HU: sources/stacks-chow/tags.txt, source lines 4580--4580. Cross-chapter standard prerequisite cited by the author inside the counted span; it is used as a pointer and its own proof is not part of this item.\par\medskip
\begin{lemma}
\label{lemma-sheaf-gives-space}
Let $S$ be a scheme and let $Y$ be an algebraic space over $S$.
Let $\mathcal{F}$ be a sheaf of sets on $Y_\etale$.
Provided a set theoretic condition is satisfied (see proof) we have
\begin{enumerate}
\item the functor $X$ associated to $\mathcal{F}$ above is an
algebraic space,
\item the map $X \to Y$ is an \'etale morphism of algebraic spaces,
\item via the identification $\Sh(Y_\etale) = \Sh(Y_{spaces, \etale})$
we have $\mathcal{F} \cong h_X$,
\item we have $\mathcal{F} \cong f_{small, !}*$. Here $*$ is the final object
of the category $\Sh(X_\etale)$ and $f_{small, !}$ exists by
Lemma \ref{lemma-etale-morphism-topoi}.
\end{enumerate}
\end{lemma}

\begin{proof}
Let us prove that $X$ is a sheaf for the fppf topology. Namely, suppose
that $\{g_i : T_i \to T\}$ is a covering of $(\Sch/S)_{fppf}$ and
$(y_i, s_i) \in X(T_i)$ satisfy the glueing condition, i.e., the
restriction of $(y_i, s_i)$ and $(y_j, s_j)$ to $T_i \times_T T_j$ agree.
Then since $Y$ is a sheaf for the fppf topology, we see that
the $y_i$ give rise to a unique morphism $y : T \to Y$ such that
$y_i = y \circ g_i$. Then we see that
$y_{i, small}^{-1}\mathcal{F} = g_{i, small}^{-1}y_{small}^{-1}\mathcal{F}$.
Hence the sections $s_i$ glue uniquely to a section of
$y_{small}^{-1}\mathcal{F}$ by
\'Etale Cohomology, Lemma \ref{etale-cohomology-lemma-describe-pullback}.

\medskip\noindent
The construction that sends $\mathcal{F} \in \Ob(\Sh(Y_\etale))$
to $X \in \Ob((\Sch/S)_{fppf})$ preserves finite limits and
all colimits since each of the functors $y_{small}^{-1}$ have
this property. Of course, if $V \in \Ob(Y_\etale)$, then
the construction sends the representable sheaf $h_V$ on $Y_\etale$
to the representable functor represented by $V$.

\medskip\noindent
By Sites, Lemma \ref{sites-lemma-sheaf-coequalizer-representable}
we can find a set $I$, for each $i \in I$ an object $V_i$ of $Y_\etale$
and a surjective map of sheaves
$$
\coprod h_{V_i} \longrightarrow \mathcal{F}
$$
on $Y_\etale$. The set theoretic condition we need is that the index
set $I$ is not too large\footnote{It suffices if the supremum of the
cardinalities of the stalks of $\mathcal{F}$ at geometric points of $Y$
is bounded by the size of some object of $(\Sch/S)_{fppf}$.}.
Then $V = \coprod V_i$ is an object of $(\Sch/S)_{fppf}$ and
therefore an object of $Y_\etale$ and we have a surjective map
$h_V \to \mathcal{F}$.

\medskip\noindent
Observe that the product of $h_V$ with itself in $\Sh(Y_\etale)$
is $h_{V \times_Y V}$. Consider the fibre product
$$
h_V \times_\mathcal{F} h_V \subset h_{V \times_Y V}
$$
There is an open subscheme $R$ of $V \times_Y V$ such that
$h_V \times_\mathcal{F} h_V = h_R$, see
Lemma \ref{lemma-support-subsheaf-final} (small detail omitted).
By the Yoneda lemma we obtain two morphisms $s, t : R \to V$ in $Y_\etale$
and we find a coequalizer diagram
$$
\xymatrix{
h_R \ar@<1ex>[r] \ar@<-1ex>[r] &
h_V \ar[r] &
\mathcal{F}
}
$$
in $\Sh(Y_\etale)$. Of course the morphisms $s, t$ are \'etale
and define an \'etale equivalence relation $(t, s) : R \to V \times_S V$.

\medskip\noindent
By the discussion in the preceding two paragraphs we find a
coequalizer diagram
$$
\xymatrix{
R \ar@<1ex>[r] \ar@<-1ex>[r] &
V \ar[r] &
X
}
$$
in $(\Sch/S)_{fppf}$. Thus $X = V/R$ is an algebraic space by
Spaces, Theorem \ref{spaces-theorem-presentation}. This proves (1).
Part (2) follows because $V \to Y$ is \'etale. Part (3) is immediate
from the definition of $X$ and $h_X$. We omit the proof of part (4);
it follows by matching the morphism associated to the cocontinuous
functor $j$ of Lemma \ref{lemma-etale-morphism-topoi}
with the description of $X_{spaces, \etale}$ as the localization
of $Y_{spaces, \etale}$ at $X$ discussed above and
Sites, Lemma \ref{sites-lemma-describe-j-shriek-representable}.
\end{proof}

\end{document}
